Large language models (LLMs) are increasingly used to label text at scale, but every label is produced by a \emph{configuration}---a model, a prompt, and a stochastic decoding run---and configurations that share training corpora and alignment procedures do not err independently. Classical aggregation models for multiple raters, from Dawid and Skene onward, assume conditional independence of raters given the true class. We show by simulation that, under the crossed dependence typical of LLM ensembles, such models increasingly overstate annotator accuracy and bias prevalence estimates as dependence grows. We propose a latent class model in which the dependence structure is derived from the factorial design of the annotation ensemble: item-level random effects that are shared across all configurations, crossed with item-by-model and item-by-prompt effects, while repeated decoding runs are treated as exchangeable. The model, estimated by Bayesian inference with the latent class marginalised, recovers point identification without a gold standard, decomposes annotation error variance into item, model, prompt and run components, and yields better-calibrated posterior class probabilities. In a simulation study it reduces the prevalence bias of Dawid--Skene several-fold (from $+0.05$ and $+0.10$ to about $+0.01$ and $+0.04$ at moderate and strong dependence), restores nominal coverage of the prevalence interval, and lowers the Brier score of the labels by a third. We study identifiability, compare the method with majority vote and Dawid--Skene, and apply it to policy-topic coding of legislative texts with human-coded reference labels. The variance decomposition translates directly into design guidance for how many models, prompts and runs to use.
